% slides/talk.tex
%
% SymbolicMOR final-project talk for 18.337.
% Three presenters, 16 minutes total: Lift, Learn, Scale.

\documentclass[aspectratio=169,10pt]{beamer}

\usetheme{Madrid}
\usecolortheme{seahorse}
\setbeamertemplate{navigation symbols}{}
\setbeamertemplate{footline}[frame number]

\usepackage{amsmath, amssymb, bm}
\usepackage{booktabs}
\usepackage{listings}
\usepackage{xcolor}

\definecolor{codebg}{RGB}{248,248,248}
\definecolor{codegreen}{RGB}{0,110,70}
\definecolor{codered}{RGB}{150,45,45}
\definecolor{codeblue}{RGB}{35,80,160}

\lstdefinelanguage{JuliaLite}{
  morekeywords={using,function,end,struct,return,for,in,if,else,elseif,try,catch,println,where,mutable,abstract,quote,let,do,module,import,export,global,local},
  sensitive=true,
  morecomment=[l]{\#},
  morestring=[b]",
}

\lstset{
  language=JuliaLite,
  basicstyle=\ttfamily\scriptsize,
  keywordstyle=\color{codeblue}\bfseries,
  commentstyle=\color{codegreen},
  stringstyle=\color{codered},
  backgroundcolor=\color{codebg},
  frame=single,
  rulecolor=\color{black!15},
  columns=fullflexible,
  keepspaces=true,
  showstringspaces=false,
  breaklines=true
}

\title[SymbolicMOR]{Symbolic Lift--Learn--Scale: \\ Intrusive Model Order Reduction in Julia}
\author{Angelo Farfan \and Prince Tufon \and Eric Noriega}
\institute{18.337 Final Project}
\date{\today}

\begin{document}

% ============================================================
% Title
% ============================================================
\begin{frame}
  \titlepage
\end{frame}

% ============================================================
% Motivation: 1 slide
% ============================================================
\begin{frame}{Motivation}
  Large nonlinear ODEs are expensive to simulate repeatedly: parameter sweeps, control loops, design optimization.

  \medskip
  \textbf{Model Order Reduction (MOR)} replaces the high-dimensional system with a small one.

  \medskip
  \textbf{Intrusive} MOR uses the equations themselves, not just data. The reduced model inherits structure (energy, conservation, stability).

  \medskip
  \textbf{Catch:} closed-form intrusive projection wants the right-hand side to be polynomial of bounded degree. Real models rarely are.

  \medskip
  \begin{block}{Our project}
    Automate the algebra. Take any nonlinear ODE, symbolically rewrite it as a quadratic system, then project. All in Julia, with parallel snapshots.
  \end{block}
\end{frame}

% ============================================================
% Pipeline overview
% ============================================================
\begin{frame}{Pipeline}
  \begin{enumerate}
    \item \textbf{Lift}: rewrite $\dot{x} = f(x)$ as $\dot{z} = G(z)$ with $G$ at most quadratic in $z$.
    \item \textbf{Learn}: collect snapshots $X = [z(t_1)\;\cdots\;z(t_K)]$, compute a POD basis $V$, project onto it.
    \item \textbf{Scale}: snapshot ensembles run in parallel; reduced-state evaluation uses sparse quadratic tensors instead of dense $H \cdot \mathrm{kron}(z,z)$.
  \end{enumerate}

  \vspace{0.4cm}
  \[
    f \;\xrightarrow{\text{symbolic}}\; G \;\xrightarrow{\text{integrate}}\; X \;\xrightarrow{\text{SVD}}\; V \;\xrightarrow{\text{Galerkin}}\; \text{ROM}
  \]

  \vspace{0.4cm}
  \begin{center}
    \footnotesize
    \textcolor{gray}{Three sections in this talk: Lift, Learn, Scale.}
  \end{center}
\end{frame}

% ============================================================
% PHASE 1: LIFT  (~5 min, ~4 slides)
% ============================================================

\begin{frame}{Phase 1: Lift. Why Polynomial Form?}
  Galerkin projection of $\dot{z} = G(z)$ onto a basis $V$ is closed-form when $G$ is a polynomial. Specifically, if
  \[
    G(z) = A z + H \, (z \otimes z) + c,
  \]
  then with $z \approx V a$,
  \[
    \dot{a} \;=\; \widehat{A}\, a + \widehat{H}\, (a \otimes a) + \widehat{c},
    \qquad
    \widehat{A} = V^{\top} A V, \quad
    \widehat{H} = V^{\top} H (V \otimes V), \quad
    \widehat{c} = V^{\top} c.
  \]

  \medskip
  No general nonlinearity in $G$ means no need to evaluate $G$ during reduced simulation. The ROM is just three small matrices.

  \medskip
  \textbf{Plan:} two symbolic rewrites that take a generic $f$ to a quadratic $G$.
  \begin{enumerate}
    \item Polynomialization (handle transcendentals, divisions).
    \item Quadratization (handle higher-degree monomials).
  \end{enumerate}
\end{frame}

\begin{frame}{Polynomialization: Chain-Rule Closure}
  For each non-polynomial atom in $f$, introduce a fresh state with closed dynamics.

  \medskip
  \begin{tabular}{ll}
    \toprule
    Atom $\phi(g)$ & New state $p$, with $\dot{p}$ \\
    \midrule
    $\exp(g)$            & $\dot{p} = p \cdot \dot{g}$ \\
    $\sin(g),\; \cos(g)$ & $\dot{p}_s = p_c \dot{g},\; \dot{p}_c = -p_s \dot{g}$ \\
    $1/g$                & $\dot{p} = -p^2 \dot{g}$ \\
    $\log(g)$            & $\dot{p} = (1/g)\, \dot{g}$ \\
    $\sqrt{g}$           & $\dot{p} = \dot{g} / (2 p)$ \\
    $\tanh(g)$           & $\dot{p} = (1 - p^2)\, \dot{g}$ \\
    \bottomrule
  \end{tabular}

  \medskip
  \textbf{Worked example.} $\dot{x} = \sin(x)$.
  Introduce $p_s = \sin(x),\; p_c = \cos(x)$:
  \[
    \dot{x} = p_s, \qquad
    \dot{p}_s = p_c\, p_s, \qquad
    \dot{p}_c = -p_s^2.
  \]
  All right-hand sides are polynomial in $(x, p_s, p_c)$. The closure is finite.
\end{frame}

\begin{frame}{Quadratization: Greedy Monomial Splitting}
  After polynomialization, every monomial of degree $> 2$ is split by introducing a fresh auxiliary equal to the product of two of its factors.

  \medskip
  \textbf{Worked example.} $\dot{x} = x^2 y^2$.

  \medskip
  Factor list: $[x, x, y, y]$. Greedy reduce:
  \begin{align*}
    [x, x, y, y] \;&\xrightarrow{\;w_1 = x \cdot x\;}\; [w_1, y, y] \\
    [w_1, y, y]  \;&\xrightarrow{\;w_2 = w_1 \cdot y\;}\; [w_2, y]
  \end{align*}
  Lifted RHS: $\dot{x} = w_2 \cdot y$.

  \medskip
  Aux dynamics by chain rule, against the post-quadratization state:
  \[
    \dot{w}_1 = 2 x \cdot \dot{x}, \qquad
    \dot{w}_2 = \dot{w}_1 \cdot y + w_1 \cdot \dot{y}.
  \]

  \medskip
  \textbf{Two correctness traps we hit:}
  \begin{itemize}\setlength\itemsep{0pt}
    \item Use the \emph{quadratized} RHS as the chain-rule basis, not the original. Otherwise high-degree monomials reappear.
    \item Reuse existing aux when the same product is requested twice. Otherwise the loop never closes.
  \end{itemize}
\end{frame}

\begin{frame}[fragile]{The \texttt{lift\_system} API}
  Takes Symbolics expressions, returns a \texttt{LiftedSystem} with original vars, lifted vars, polynomial RHS $G$, and aux definitions.

\begin{lstlisting}
using Symbolics, SymbolicMOR
@variables x

ls = lift_system([x], [tanh(x)])
# ls.original_vars : [x]
# ls.lifted_vars   : [x, p1, w1]
# ls.aux_eqs       : [p1 ~ tanh(x), w1 ~ p1^2]
# ls.F             : [p1, p1 - p1*w1, 2*w1 - 2*w1^2]
\end{lstlisting}

  \medskip
  Verified end to end against the original ODE for
  $\exp,\;\sin,\;\cos,\;1/g,\;\log,\;\sqrt{\;\;},\;\tanh$. Lifted trajectories match unlifted ones to $\sim 10^{-10}$.

  \medskip
  Also accepts a ModelingToolkit \texttt{ODESystem} via \texttt{lift\_system(sys)}.
\end{frame}

% ============================================================
% PHASE 2: LEARN  (~5 min, ~5 slides)
% ============================================================

\begin{frame}{Phase 2: Learn. POD Basis from Snapshots}
  Collect $K$ trajectory snapshots of the lifted system into a matrix
  \[
    X = \begin{bmatrix} z(t_1) & z(t_2) & \cdots & z(t_K) \end{bmatrix} \in \mathbb{R}^{N \times K}.
  \]

  \medskip
  \textbf{Mean-center} before SVD. Compute the average snapshot $\bar{z}$ and subtract:
  \[
    \bar{z} \;=\; \frac{1}{K}\sum_{j=1}^{K} z(t_j),
    \qquad
    \widetilde{X} \;=\; X - \bar{z}\,\mathbf{1}^{\top}.
  \]
  Then $\widetilde{X} = U\Sigma W^{\top}$, and $V = U[:, 1{:}r] \in \mathbb{R}^{N \times r}$ is the optimal rank-$r$ basis for the centered snapshots.

  \medskip
  Reconstruction of the lifted state and the reduced coordinates:
  \[
    z \;\approx\; \bar{z} + V a,
    \qquad
    a \;=\; V^{\top}\,(z - \bar{z}).
  \]
  $V$ captures deviations around the mean, which usually gives a better low-rank fit than uncentered POD. (Set $\bar{z} = 0$ to recover the uncentered formulation.)

  \medskip
  Energy retention: choose $r$ so that $\sum_{i \leq r} \sigma_i^2 \,/\, \sum_i \sigma_i^2 \geq $ a threshold (e.g. $99.99\%$).
\end{frame}

\begin{frame}{Two ROM Forms (We Implemented Both)}
  Once we have $V$ and the lifted system $\dot{z} = G(z)$, the reduced model can be written two ways.

  \begin{columns}[T]
    \begin{column}{0.49\textwidth}
      \textbf{Simple (function-evaluating)}
      \[
        \dot{a} \;=\; V^{\top}\, G(\bar{z} + V a)
      \]

      \medskip
      Pedagogically the cleanest. One line of code given $V$ and a callable $G$.

      \medskip
      Calls the full lifted RHS at every step of the reduced ODE: $O(N)$ work per step.
    \end{column}
    \begin{column}{0.49\textwidth}
      \textbf{Explicit (operator-precomputed)}
      \[
        \dot{a} \;=\; \widehat{A}\, a + \widehat{H}\, (a \otimes a) + \widehat{c}
      \]

      \medskip
      \[
        \widehat{A} = V^{\top} A V,\;\; \widehat{H} = V^{\top} H (V \otimes V),\;\; \widehat{c} = V^{\top} c
      \]

      \medskip
      Pre-compute once, then $O(r^2)$ work per step. Independent of $N$ at runtime.
    \end{column}
  \end{columns}

  \medskip
  Both produce identical $\dot{a}$ in our tests. The explicit form is the production path; the simple form is what we open the live demo with.
\end{frame}

\begin{frame}{Why $\widehat{H}$ Hurts: the $\mathrm{kron}(V, V)$ Wall}
  The naive formula has a serious cost.
  \[
    \widehat{H} \;=\; V^{\top}\, H\, (V \otimes V).
  \]

  \medskip
  Sizes: $V$ is $N \times r$, so $V \otimes V$ is $N^2 \times r^2$. For $N = 200$ and $r = 8$, that is a $40000 \times 64$ matrix. For $N = 1000$, it is $10^6$ rows.

  \medskip
  Forming $V \otimes V$ explicitly is the bottleneck, not the multiply.

  \medskip
  \begin{block}{Our solution: \texttt{QuadraticTensor}}
    Store $H$ as $(i, p, q, \text{val})$ tuples, then iterate the contraction
    \[
      \widehat{H}_{\alpha \beta \gamma} \;=\; \sum_{(i, p, q)} V_{i \alpha}\, V_{p \beta}\, V_{q \gamma}\, H_{ipq}
    \]
    skipping zeros. No $\mathrm{kron}(V, V)$ ever materialized.
  \end{block}
\end{frame}

\begin{frame}{Sparse Tensor Galerkin: One Diagram}
  \[
    H_{ipq}\;\;\text{nonzero only for a few}\;(p,q)\;\text{per}\;i.
  \]

  \begin{center}
    \begin{tabular}{lcc}
      \toprule
      Representation             & Bytes for $\widehat{H}$ & Build time \\
      \midrule
      Dense $H \cdot (V \otimes V)$  & $O(r \cdot N^2)$  &  large \\
      Sparse $H$, dense $V \otimes V$ & $O(r \cdot N^2)$  &  large \\
      \texttt{QuadraticTensor}        & $O(\mathrm{nnz}(H) \cdot r^3)$ & moderate \\
      \bottomrule
    \end{tabular}
  \end{center}

  \medskip
  In our benchmarks the tensor path is the only one that scales when $N$ exceeds a few hundred. All three paths agree numerically: a verified part of our test suite.

  \medskip
  Implementation lives in \texttt{src/learn/quadratic\_operator.jl} and \texttt{pod\_galerkin.jl}.
\end{frame}

\begin{frame}[fragile]{End-to-End Code}
\begin{lstlisting}
using SymbolicMOR, Symbolics, OrdinaryDiffEq

@variables x y z
rhs = [10.0*(y-x), x*(28.0-z) - y, x*y - (8/3)*z]
ls  = lift_system([x, y, z], rhs)            # Phase 1

# Snapshots from many ICs (Phase 3 will parallelize)
G! = build_lifted_rhs(ls)
u0s = [randn(3) .+ [1.0, 0.0, 25.0] for _ in 1:20]
X   = generate_snapshots(G!, u0s, (0.0, 5.0))

V, _, _ = compute_pod_basis(X, 4)             # Phase 2
A, Q, c = extract_quadratic_tensor(ls)
A_hat, Q_hat, c_hat = galerkin_project(A, Q, c, V)

# Reduced ODE: 3 small operators, no full RHS calls
prob = ODEProblem(rom_rhs!, V'*u0s[1], (0.0, 5.0),
                  (A_hat, Q_hat, c_hat))
\end{lstlisting}
\end{frame}

% ============================================================
% PHASE 3: SCALE  (~3 min, ~3 slides)
% ============================================================

\begin{frame}{Phase 3: Scale. Snapshot Ensembles in Parallel}
  Snapshot generation is embarrassingly parallel: each initial condition runs an independent ODE solve. We use \texttt{Distributed} so trajectories spread across workers (and across cluster nodes if available).

  \medskip
  \begin{itemize}
    \item \texttt{generate\_snapshots\_parallel}: each worker integrates a disjoint subset of ICs.
    \item \texttt{generate\_snapshots\_ensemble}: SciML \texttt{EnsembleProblem} backend, integrates with \texttt{EnsembleThreads} or \texttt{EnsembleDistributed}.
    \item Both return a stacked snapshot matrix that feeds straight into the POD step.
  \end{itemize}

  \medskip
  Per-trajectory cost dominates over reduction in serial; once parallel, reduction becomes the bottleneck. That motivated the tensor work in Phase 2.
\end{frame}

\begin{frame}{Validation: Lift Accuracy and ROM Convergence}
  \begin{columns}[T]
    \begin{column}{0.52\textwidth}
      \textbf{Lift accuracy.} Max trajectory error of the lifted system against the unlifted one over the test horizon, with \texttt{Tsit5} at \texttt{abstol = reltol = 1e-10}.

      \medskip
      \footnotesize
      \begin{tabular}{lccc}
        \toprule
        System                  & $N$ lifted & Aux & Lift err \\
        \midrule
        Cubic decay $-x^3$      & 2 & 1 & $1.5{\cdot}10^{-10}$ \\
        Van der Pol             & 3 & 1 & $3.2{\cdot}10^{-10}$ \\
        Lorenz (already poly)   & 3 & 0 & $0$ \\
        Allen--Cahn ($-u^3$)    & 2 & 1 & $1.2{\cdot}10^{-10}$ \\
        \bottomrule
      \end{tabular}
    \end{column}
    \begin{column}{0.45\textwidth}
      \textbf{ROM convergence on Lorenz.} Energy retained vs reduced-state ROM error over $t \in [0, 2]$, $K = 9624$ snapshots from $24$ initial conditions.

      \medskip
      \footnotesize
      \begin{tabular}{ccc}
        \toprule
        $r$ & energy & ROM err \\
        \midrule
        1 & $0.82$ & $4.3{\cdot}10^{1}$ \\
        2 & $0.99$ & $3.6{\cdot}10^{1}$ \\
        3 & $1.00$ & $1.4{\cdot}10^{-9}$ \\
        \bottomrule
      \end{tabular}

      \medskip
      \scriptsize Lorenz is chaotic, so $r < 3$ ROMs lose phase quickly even with high energy retention. $r = 3$ recovers the full system to solver tolerance.
    \end{column}
  \end{columns}

  \medskip
  Test suite: 596 tests, runs in $\sim$3 s after package load. CI on every push.
\end{frame}

\begin{frame}{Tensor vs Dense Quadratic Evaluation}
  Evaluate $H \, (z \otimes z)$ at lifted dimension $n$, with random sparsity $\mathrm{nnz}(H) = 4n$ (a regime typical of discretized PDEs after lifting).

  \medskip
  \begin{center}
    \footnotesize
    \begin{tabular}{rcccc}
      \toprule
      $n$ & dense $H \cdot \mathrm{kron}$ (ms) & \texttt{QuadraticTensor} (ms) & speedup & memory ratio \\
      \midrule
      40   & 0.004  & < 0.001 & 20x   & 100x   \\
      80   & 0.066  & < 0.001 & 150x  & 400x   \\
      160  & 0.734  & 0.001   & 780x  & 1600x  \\
      320  & 7.55   & 0.002   & 4170x & 6400x  \\
      \bottomrule
    \end{tabular}
  \end{center}

  \medskip
  Dense scales as $O(n^3)$. The tensor evaluation visits only the stored entries, so it stays roughly proportional to $\mathrm{nnz}(H)$ regardless of $n$. The same crossover applies to the \emph{reduced} operator $\widehat{H}$ once $r$ is large enough.

  \medskip
  Reproduce: \texttt{julia --project=. benchmarks/quadratic\_operator\_benchmark.jl}
\end{frame}

% ============================================================
% Wrap-up
% ============================================================

\begin{frame}{What Worked, What Bit Us, What Is Next}
  \textbf{Worked}
  \begin{itemize}\setlength\itemsep{0pt}
    \item Symbolics + SymbolicUtils made the algebra tractable. The chain-rule closures for the seven supported atoms are short.
    \item Two ROM forms in the same library: a clear pedagogical contrast and a practical fast path.
    \item Tensor representation kept the projection feasible for nontrivial $N$.
  \end{itemize}

  \medskip
  \textbf{Bit us}
  \begin{itemize}\setlength\itemsep{0pt}
    \item Pattern-matching substitution misses subterms inside multi-arg products. Required expanding before substituting and deduplicating aux equations.
    \item Numeric literals inside expanded expressions are wrapped (not bare \texttt{Number}) in SymbolicUtils 4.x. Custom polynomial-detection had to unwrap.
  \end{itemize}

  \medskip
  \textbf{Next}
  \begin{itemize}\setlength\itemsep{0pt}
    \item More general non-polynomial atoms (rational, piecewise).
    \item Automatic basis selection beyond fixed energy thresholds.
    \item GPU kernels for the tensor contraction.
  \end{itemize}
\end{frame}

\begin{frame}{References and Code}
  \footnotesize
  \begin{itemize}
    \item Kramer \& Willcox (2019). \emph{Nonlinear Model Reduction via Lift \& Learn}. AIAA Journal.
    \item Bychkov \& Pogudin (2021). \emph{Optimal Monomial Quadratization for ODE Systems}. arXiv:2103.08013.
    \item Benner, Gugercin \& Willcox (2015). \emph{Survey of Projection-Based Model Reduction}. SIAM Review.
  \end{itemize}

  \medskip
  \textbf{SymbolicMOR.jl}
  \begin{itemize}
    \item Repo: \texttt{github.com/elteomc/SymbolicMOR}, branch \texttt{quad-tensor-ops}.
    \item Built on Symbolics.jl, ModelingToolkit.jl, OrdinaryDiffEq.jl, Distributed.
    \item Documentation: \texttt{docs/} (Documenter.jl), live walkthrough in \texttt{scripts/lorenz\_demo.jl}.
  \end{itemize}

  \vspace{0.3cm}
  \begin{center}
    \large Questions?
  \end{center}
\end{frame}

% ============================================================
% Backup slides (appendix). Not counted toward the 16-minute budget.
% Use them in Q&A.
% ============================================================

\appendix

\begin{frame}{Backup: Which Functions Can We Lift?}
  Lifting an atom $\phi(g)$ requires a finite chain-rule closure. We pick a small fresh state $p_1, p_2, \ldots$ such that each $\dot{p}_i$ is polynomial in the existing states.

  \medskip
  \begin{columns}[T]
    \begin{column}{0.48\textwidth}
      \textbf{Closes finitely (supported)}
      \begin{itemize}\setlength\itemsep{0pt}
        \item $\exp(g)$. Self-closing.
        \item $\sin(g),\;\cos(g)$. Close as a pair.
        \item $1/g,\;g^{-k}$. Self-closing.
        \item $\log(g)$. Closes via $1/g$.
        \item $\sqrt{g}$. Closes via $1/p$.
        \item $\tanh(g)$. Self-closing via $(1 - p^2)$.
      \end{itemize}
    \end{column}
    \begin{column}{0.48\textwidth}
      \textbf{Does not close finitely}
      \begin{itemize}\setlength\itemsep{0pt}
        \item $|x|$. Not differentiable at zero.
        \item Piecewise functions. Need event handling, not algebraic closure.
        \item $\arctan(g),\;\arcsin(g)$. Possible in principle (each closes through a rational extension), not yet implemented.
        \item Arbitrary user functions. No general procedure.
      \end{itemize}
    \end{column}
  \end{columns}

  \medskip
  Adding a new closing atom is one rule in \texttt{src/lift/polynomialize.jl}: name the atom, give the chain-rule expression, register sibling atoms if needed.
\end{frame}

\begin{frame}{Backup: Greedy vs QBee Quadratization}
  \begin{columns}[T]
    \begin{column}{0.48\textwidth}
      \textbf{What we use (greedy)}
      \begin{itemize}\setlength\itemsep{0pt}
        \item Locate a monomial of degree $> 2$.
        \item Factor into variables, take the first two factors.
        \item Replace them with a fresh aux $w_k = f_1 \cdot f_2$.
        \item Repeat. Apply chain rule for each new $w_k$.
      \end{itemize}
      Simple, predictable, fast on the systems we tested. Can introduce more aux variables than strictly necessary.
    \end{column}
    \begin{column}{0.48\textwidth}
      \textbf{QBee-style search (Bychkov \& Pogudin, 2021)}
      \begin{itemize}\setlength\itemsep{0pt}
        \item Searches the monomial structure globally.
        \item Looks for the smallest set of aux monomials that closes the system.
        \item Solved as integer programming over an explicit model.
        \item Optimal in lifted dimension.
      \end{itemize}
      Heavier engineering. Wins when greedy gets stuck or balloons.
    \end{column}
  \end{columns}

  \medskip
  We did not implement QBee in the time we had. The dedup step in our quadratization (reuse an aux when the same product is requested twice) recovers some of QBee's economy in practice.
\end{frame}

\end{document}
